% !PRINTED 618
\weberpage{618}
%<<CONT
is not a double root. For the different roots of this resolvent belong to conjugate groups, and if therefore one of these roots is rational, then one of the conjugate groups is metacyclic. We therefore have the theorem:

\gesperrt{V. The necessary and sufficient condition for the solvability of the equation $f(x)=0$ by radicals is that the resolvent of degree $\nu$, $F(y)=0$, which one has arranged, by a suitable choice of the function $y$, so that it receives no double roots, has a rational root.}

Another form of this condition results in the following way.

Among the permutations of a linear group, only the identical one leaves any two figures unchanged. For if $(z, az+b)$ does not alter two figures, then the congruence $az+b \equiv z \pmod n$ must have two distinct solutions. But this is possible only if $a \equiv 1$, $b \equiv 0 \pmod n$.

If, then, the group $P$ of $f(x)$ is metacyclic, it reduces itself, by the adjunction of two arbitrary roots, to the unit, and consequently all the roots are rationally expressible in terms of any two among them. Hence:

\gesperrt{VI. If an irreducible equation of prime degree is metacyclic, then all the roots are rationally expressible in terms of any two among them.}

But this theorem also holds conversely, which we can prove simply in the following manner.

Suppose an irreducible equation of prime degree $n$ has the property that all the roots are rationally expressible in terms of any two among them. If, say, $x_\nu = \psi(x_0, x_1)$ is such a representation, then all the permutations of the group $P$ of our equation can be applied to this relation, and if therefore one of these permutations leaves $x_0$ and $x_1$ unchanged, then it also leaves all the $x_\nu$ unchanged, i.e. it is the identical permutation. \gesperrt{Hence $P$ contains, besides the identical permutation, none that leaves two figures unaltered.}

If now one of the permutations $\pi = cc_1 \ldots$ of $P$ contains two or more different cycles $c, c_1 \ldots$, and $c$ is of degree $h$, but $c_1$ of a higher degree, then $\pi^h = c_1 \ldots$, and $\pi^h$ leaves
%CONT>>
